% ============================================================================
% Drop-in replacement for the paper's mathematical core (new Sections 2-3).
% Requires theorem environments: theorem, proposition, lemma, corollary,
% conjecture, remark (numbered within section), and \R, \M macros:
%   \newcommand{\R}{\mathbb{R}}  \newcommand{\M}{\mathrm{M}}
% All proofs are in appendix_proofs.tex (\appendixproofs labels).
% Every numbered statement below was verified numerically (float64, machine
% precision) by two independent adversarial verification agents; see the
% reproducibility note in the appendix.
% ============================================================================

\section{The knowledge matrix is the germ of the network function}
\label{sec:germ}

Throughout, $f(\theta,\cdot):\R^d\to\R^C$ is a feedforward network with
piecewise-linear (PL) components at inference time --- ReLU-family activations,
affine layers (convolutions, batch normalisation in evaluation mode), and
max/average pooling; this covers every architecture in our experiments. Write
$D^{(\ell)}(x)$ for the $0/1$ activation masks at input $x$,
$J(\theta,x)=W^{(L)}D^{(L-1)}(x)\cdots D^{(1)}(x)W^{(1)}$ for the masked-product
Jacobian, and $c(\theta,x)$ for the accumulated bias; on the activation region
of $x$ the network is the affine map $x'\mapsto Jx'+c$. The knowledge matrix is
\begin{equation}
\M(\theta,x)\;=\;\bigl[\,J(\theta,x)\,\mathrm{diag}(x)\;\big|\;c(\theta,x)\,\bigr]
\;\in\;\R^{C\times(d+1)},
\qquad
\M(\theta,x)\,\mathbf 1_{d+1}\;=\;f(\theta,x)
\quad\text{(exact, at every $x$)} .
\label{eq:km-def}
\end{equation}

We first fix the object that the knowledge matrix records. Call
$X_{\mathrm{reg}}\subseteq\R^d$ the \emph{regular set} of inputs whose
activation pattern is locally constant; it is open, dense, and of full
Lebesgue measure (Lemma~\ref{lem:null}), and on it the network coincides with a
single affine map on a neighbourhood of each point.

\begin{definition}[Germ of a PL network]\label{def:germ}
Let $f(\theta,\cdot):\R^d\to\R^C$ be piecewise linear and let $x\in
X_{\mathrm{reg}}$. Two functions agreeing on \emph{some} neighbourhood of $x$
define the same \emph{germ} at $x$; the germ of $f$ at $x$ is the equivalence
class of $f$ under this relation. Because $f$ is affine on a neighbourhood of
$x$, this class is represented by the local affine pair
$\bigl(J(\theta,x),\,c(\theta,x)\bigr)$ with $f(x')=J(\theta,x)\,x'+c(\theta,x)$
for all $x'$ near $x$; equivalently, by the first-order Taylor data
$\bigl(Df(x),\,f(x)\bigr)$ since $J(\theta,x)=Df(x)$ and $c(\theta,x)=f(x)-Df(x)\,x$.
The germ depends on $\theta$ only through the function $f(\theta,\cdot)$ realises
locally, not through how $\theta$ encodes it. Throughout, ``germ'' means this
local affine pair.
\end{definition}

\begin{theorem}[Germ identity]\label{thm:germ}
Let $X_{\mathrm{reg}}\subseteq\R^d$ be the (open, dense, full-measure) set of
inputs whose activation pattern is locally constant. For every $x\in
X_{\mathrm{reg}}$, $f(\theta,\cdot)$ is affine on a neighbourhood of $x$ and
\[
J(\theta,x)=Df(x),\qquad c(\theta,x)=f(x)-Df(x)\,x,\qquad
\M(\theta,x)=\bigl[\,Df(x)\,\mathrm{diag}(x)\;\big|\;f(x)-Df(x)\,x\,\bigr].
\]
In particular $\M(\theta,x)$ depends on $\theta$ only through the germ of the
realised function at $x$. At boundary points the identity
$\M\mathbf 1=f(x)$ still holds exactly, and the mask convention
$\mathbb 1[z>0]$ selects the one-sided germ from the inactive side.
\end{theorem}

\begin{remark}[Where the columns vanish]\label{rem:zero-pixels}
Column $i$ of $\M(x)$ is $J_{:,i}\,x_i$: on $\{x_i=0\}$ it vanishes
identically, so the $i$-th germ column is not recoverable from $\M(x)$ there.
This set is Lebesgue-null but has \emph{positive probability} under raw image
data; the condition must be checked on the preprocessed input the network
sees. The variant $W_{\mathrm{eff}}(x)=[\,J\mid c\,]$ recovers the germ with no
condition on $x$. Invariance (Theorem~\ref{thm:maxinv}) needs no such
hypothesis; germ \emph{recovery} (Theorem~\ref{thm:complete}(i)) does.
\end{remark}

\subsection{Invariance: the largest possible group}
\label{sec:invariance}

The knowledge matrix is built functorially from the quiver representation of
the network \citep{armenta2021representation,armenta2022double}, and this fixes
the invariance picture before any analysis. Two sources of invariance must be
kept apart, because they need entirely different machinery.

\emph{Quiver isomorphisms} act on a network without changing the underlying
quiver representation up to isomorphism: relabelling (permuting) the neurons in
a hidden layer and rescaling them by nonzero per-neuron factors. Each induces
an isomorphism of quiver representations, and the knowledge matrix --- being a
functorial construction on that representation --- is invariant by
construction. This is one telescoping lemma, stated next; positive rescalings
(the ReLU case of Thms.~4.1--4.2 of the original version, after
\citealp{armenta2021representation}) and neuron permutations are \emph{both}
elements of this group and are covered by it identically. We do \emph{not}
attribute permutation invariance to Thms.~4.1--4.2: those theorems are about
rescalings, and permutation invariance is a separate special case of the same
quiver-isomorphism principle, not a corollary of them.

\emph{Function-preserving changes that are not quiver isomorphisms} --- dead-unit
insertion, neuron splitting, and above all moving to a \emph{different
architecture} that realises the same function near $x$ --- lie in the full
stabiliser of the germ but outside the quiver-isomorphism group. Invariance
under these is exactly what is \emph{not} automatic, and it is the genuinely
novel content of the germ identity (Theorem~\ref{thm:maxinv}): the matrix
depends on $\theta$ only through the germ of Definition~\ref{def:germ}, so any
two encodings of the same local function --- however different their graphs ---
produce the same matrix.

\begin{lemma}[Quiver-isomorphism invariance]\label{lem:quiver-inv}
Let $\theta'$ be obtained from $\theta$ by a quiver isomorphism: a hidden-layer
neuron permutation, a nonzero per-neuron rescaling, or any composite thereof,
with the per-neuron activations carried with their units. Then
$\M(\theta',x)=\M(\theta,x)$ for every $x\in\R^d$ at which both networks have a
well-defined secant decomposition, by the conjugation/telescoping of the slope
products and bias accumulators (proof of Proposition~\ref{prop:perm}; positive
rescalings telescope by the same argument).
\end{lemma}

\begin{theorem}[Maximal invariance: across architectures]\label{thm:maxinv}
Let $N_1,N_2$ be any two PL networks --- any architectures, widths, depths,
parameters --- realising the same function on a neighbourhood of
$x\in X_{\mathrm{reg}}(N_1)\cap X_{\mathrm{reg}}(N_2)$. Then
$\M_1(x)=\M_2(x)$. Consequently $\M(\cdot)(x)$ is invariant under the
\emph{entire stabiliser} of the realised germ (Definition~\ref{def:germ}) ---
the largest group under which any function-level observable can possibly be
invariant. The cross-architecture case is the proper payload of this theorem:
it is \emph{not} a quiver isomorphism (Lemma~\ref{lem:quiver-inv}) and could
not follow from it.
\end{theorem}

\begin{proposition}[Permutation invariance, every input]\label{prop:perm}
Neuron permutation is a quiver isomorphism, so $\M$-invariance under it is
immediate from Lemma~\ref{lem:quiver-inv}; the only non-automatic content is its
\emph{scope}. The telescoping needs elementwise (neuron-wise equivariant)
activations, so that the permutation matrix $P$ commutes with the nonlinearity
$\varphi$; non-equivariant vector activations are excluded. Within that scope,
for hidden-layer permutations (weights conjugated, neuron-wise activations
carried with the units), $\M(\theta^\pi,x)=\M(\theta,x)$ for \emph{every}
$x\in\R^d$, including region boundaries, under the secant convention. The scope
is sharp: for activations with $\varphi(0)\neq0$ the row-sum identity itself can
fail on the measure-zero set $\{z=0\}$ (where $P$ no longer commutes with the
constant offset $\varphi(0)$).
\end{proposition}

So permutations and positive rescalings telescope identically; neither is the
``gap'' --- both are quiver isomorphisms. \emph{Function-preserving surgery}
(dead-unit insertion, neuron splitting, architecture change) is the part that
genuinely needs the germ identity, and is covered at generic $x$ by
Theorem~\ref{thm:maxinv}. Genericity is not removable: the
identity map and $g(x)=\mathrm{ReLU}(x-1)-\mathrm{ReLU}(1-x)+1\equiv x$ have
$\M_{\mathrm{id}}(1)=[\,1\mid 0\,]\neq[\,0\mid 1\,]=\M_g(1)$ at the kink (equal
row sums). Conversely, the stabiliser is strictly larger than the isomorphism
group (Proposition~\ref{prop:contraction}), so $\M$ cannot determine a network
up to isomorphism --- a feature for robustness claims, a boundary for inverse
problems.

\subsection{Completeness: the title claim, as a theorem}
\label{sec:completeness}

\begin{theorem}[Hidden activations are not enough]\label{thm:complete}
\emph{(i)} For $x\in X_{\mathrm{reg}}$ with all $x_i\neq0$, $\M(x)$ determines
the germ $(J,c)$ of $f$ at $x$. \emph{(ii)} The field $x\mapsto\M(x)$
determines $f$ exactly on any set ($f=\M\mathbf 1$ pointwise), and globally
when the set has full measure; on a mere open set it does \emph{not} determine
$f$ beyond the closures of the regions it meets (sharp: $\mathrm{ReLU}(x{-}1)$
and $\mathrm{ReLU}(x)$ have identical fields on $(-1,0)$).
\emph{(iii)} Hidden activations are gauge-covariant: under a positive
per-neuron rescaling, $a_v\mapsto\tau_v a_v$ while the function is unchanged
--- any activation statistic that is not gauge-invariant is not a function of
$(f,x)$. \emph{(iv)} Hidden activations are germ-incomplete: for any $x$ with
strict masks, $Df(x)\neq0$, and at least one hidden layer, there are
arbitrarily small parameter changes that leave \emph{every} hidden pre- and
post-activation at $x$ and the output $f(x)$ unchanged while changing the germ
--- hence changing $f$ on every neighbourhood of $x$, and changing $\M(x)$.
\end{theorem}

Parts (iii) and (iv) together are the precise form of the title: activations
are neither invariant (two networks, same function, different activations
everywhere) nor complete (two networks, identical activations at $x$,
different functions near $x$). The pair $(D(x),\M(x))$ is both --- a complete,
gauge-invariant record of the computation's germ, with the completeness
hypothesis ($x_i\neq0$) stated by Remark~\ref{rem:zero-pixels}.

\subsection{Relation to attribution methods}
\label{sec:attribution}

\begin{theorem}[Knowledge matrix $=$ gradient$\times$input $\oplus$ bias attribution]\label{thm:weff}
For PL networks at $x\in X_{\mathrm{reg}}$ (no pooling ties), the first $d$
columns of $\M(x)$ are exactly per-class gradient$\times$input,
$\M_{:,i}(x)=(\partial f/\partial x_i)\,x_i$, and the last column is the raw,
signed, class-resolved aggregate of the network's bias attributions,
$c=\sum_\ell(\partial f/\partial b^{(\ell)})\,b^{(\ell)}$. Hence $\M(x)$ equals
per-class gradient$\times$input augmented by an exact aggregate bias column. The
same matrix is computed by several equivalent routes (masked product, the probe
construction, and autograd), detailed in Appendix~\ref{app:engineering}.
\end{theorem}

\textbf{We own this equivalence as a strength.} For PL networks at almost every
input, the knowledge matrix coincides, column for column, with per-class
gradient$\times$input \citep{shrikumar2016not,ancona2018towards} augmented by a
per-class aggregate bias attribution in the sense of FullGrad
\citep{srinivas2019full}; the per-region affine operator itself is the object of
the spline view of deep networks \citep{balestriero2018spline} and of Jacobian
analyses of bias-free denoisers \citep{mohan2020robust}. We therefore flag, as a
load-bearing premise of everything below: \emph{every function-level property of
$\M$ on PL nets --- invariance, completeness, fixed shape, exact row sums --- is
shared by that gradient$\times$input data, not specific to the knowledge
matrix.} This is not a concession but a foothold: it grounds the construction in
established attribution theory, supplies its exact completeness identity, and is
what makes the vector--Jacobian-product computation of Appendix~\ref{app:engineering}
possible. What the knowledge matrix adds over the same data is not a new
pointwise invariant but a \emph{canonical arrangement}: one fixed-shape matrix
per input with exact row-sum accounting (the basis of
Section~\ref{sec:geometry}), comparable across architectures without neuron
alignment, with a path-resolved quiver-representation lift above it
\citep{armenta2021representation,armenta2022double} and a secant extension to
non-PL activations where gradient attributions satisfy no exact completeness
identity. Computation of $\M$ --- the probe construction, the $C$
vector--Jacobian-product route, the measured speed-up, and numerical validation
--- is deferred to Appendix~\ref{app:engineering}.

\begin{remark}[The quiver lift is scaffolding for future work, not machinery of this paper]
\label{rem:series-positioning}
Every result in this paper --- the germ identity (Theorem~\ref{thm:germ}), maximal
invariance (Theorem~\ref{thm:maxinv}), completeness (Theorem~\ref{thm:complete}),
and the entire distance geometry of Section~\ref{sec:geometry} --- follows from
the germ alone, i.e.\ from the chain rule on the active region. We use the
path-resolved quiver representation nowhere in this paper; the knowledge matrix is a
\emph{contraction} of it (Proposition~\ref{prop:contraction}). We state this
plainly rather than leave it as an invited objection: the quiver representation
strictly refines the matrix --- it separates functionally identical but
non-isomorphic networks that the matrix identifies --- and it is the object of
future work where that finer information is needed. In this paper it
is scaffolding for the construction, not a load-bearing tool.
\end{remark}

\begin{proposition}[The per-input contraction is exactly what is lost]
\label{prop:contraction}
There exist same-architecture ReLU networks, functionally identical --- hence
with $\M_A(x)=\M_B(x)$ for every $x$ --- that are non-isomorphic, with
distinct path-value multisets; the path-resolved quiver representation
distinguishes them, the matrix cannot. On the identifiable architecture class
of \citet{phuong2020functional} (and within general networks of the same
architecture), knowledge-matrix-field equality is equivalent to
permutation$\times$positive-rescaling equivalence; for general architectures
the fibre question is open and connects to functional dimension
\citep{grigsby2022functional} and the neuromanifold fibres of
\citet{flinth2026fibers}.
\end{proposition}

